\subsection{\texorpdfstring{S-bornologies on \(\mathcal{L}(\mathbf{E};\mathbf{F})\)}{S-bornologies on \textbackslash mathcal\{L\}(\textbackslash mathbf\{E\};\textbackslash mathbf\{F\})}}

Let E and F be two locally convex spaces and S a family of bounded subsets of E. To say that a subset H of \(\mathcal{L}(E; F)\) is bounded for the S-topology means that for every \(\mathbf{M} \in \mathfrak{S}\) , every neighbourhood \(V\) of 0 in \(F\) absorbs the set \(\mathrm{H(M)} = \bigcup_{u \in H} u(\mathrm{M})\) ; this is the same as saying that for every \(\mathbf{M} \in \mathfrak{S}\) , the set \(\mathrm{H(M)}\) is bounded in \(F\) . Equivalently, this means that for every neighbourhood \(V\) of 0 in \(F\) , the set \(\bigcap_{u \in H} u^{-1}(V)\)

absorbs every subset M of S.

PROPOSITION 9. --- Let E and F be two locally convex spaces and S a family of bounded subsets of E. Then every equicontinuous subset of \(\mathcal{L}(\mathrm{E};\mathrm{F})\) is bounded for the S-topology.

For, if \(\mathbf{H}\) is an equicontinuous subset of \(\mathcal{L}(\mathrm{E};\mathrm{F})\) and \(\mathbf{V}\) a neighbourhood of 0 in \(\mathbf{F}\) , the set \(\bigcap_{u\in \mathbf{H}}u^{-1}(\mathbf{V})\) is a neighbourhood of 0 in \(\mathbf{E}\) , hence absorbs every bounded subset of \(\mathbf{E}\) .

A subset of \(\mathcal{L}(\mathrm{E};\mathrm{F})\) which is bounded for a S-topology is not necessarily equicontinuous, even if S is covering and S is the canonical bornology on E (IV, p.~50, exerc. 17). In the following paragraph we shall study, under the name barrelled spaces, the spaces E such that every simply bounded subset of \(\mathcal{L}(\mathrm{E};\mathrm{F})\) is equicontinuous. For the present note the following result:

PROPOSITION 10. --- Let E be a bornological space (in particular, a metrizable locally convex space) and F a locally convex space. Every subset H of \(\mathcal{L}(\mathrm{E};\mathrm{F})\) which is bounded for the topology of bounded convergence is equicontinuous.

For every convex balanced neighbourhood V of 0 in F, the set \(\bigcap_{u\in\mathbf{H}}u^{-1}(\mathbf{V})\) absorbs every bounded subset of E, hence is a neighbourhood of 0 in E; this proves that H is equicontinuous.

